\chapter{Entropy, Codes and Typical Sequences}

\begin{center}\itshape
A fair coin has entropy $1$ bit. A coin that lands heads with probability $0.1$\\
has entropy about $0.469$ bits. The typical sequences are the ones you actually see.
\end{center}

Concentration says a sum settles near its mean. Entropy says which messages are typical. The two laws answer different questions. One is about the size of a fluctuation. The other is about the size of the set of strings that carry the probability.

\section{Entropy of a finite alphabet}

Let $X$ take finitely many values, with $P(X=a)=p_a>0$. The \emph{Shannon entropy} in bits is
\[
 H(X)=\sum_a p_a\log_2\frac1{p_a}.
\]
The sum is finite. Changing the base of the logarithm multiplies $H$ by a constant; bits are the choice $\log_2$. A fair coin has $H=1$. A coin with $P(\mathrm{heads})=0.1$ has
\[
 H=0.1\log_2 10+0.9\log_2(10/9)=0.46899\ldots.
\]

\begin{proposition}[Gibbs]\label{prop:gibbs}
If $X$ takes at most $m$ values, then $H(X)\le\log_2 m$, with equality if and only if $X$ is uniform on $m$ values.
\end{proposition}
\begin{proof}
The inequality $\ln u\le u-1$ for $u>0$, from Volume~1, Chapter~9, gives, for any probability masses $q_a>0$ on the same values,
\[
 \sum_a p_a\ln\frac{q_a}{p_a}\le\sum_a p_a\Bigl(\frac{q_a}{p_a}-1\Bigr)=\sum_a q_a-\sum_a p_a=0.
\]
Thus $\sum p_a\log_2(1/p_a)\le\sum p_a\log_2(1/q_a)$. Take $q_a=1/m$. Equality holds when $q_a=p_a$ for every $a$, because $\ln u=u-1$ only at $u=1$.
\end{proof}

For i.i.d.\ copies $X_1,X_2,\ldots$ with the same finite law, the probability of a particular string is the product of the symbol probabilities, so
\[
 -\frac1n\log_2 P(X_1=a_1,\ldots,X_n=a_n)=\frac1n\sum_{i=1}^n\log_2\frac1{p_{a_i}}.
\]
The summands are i.i.d.\ and bounded, because every $p_a$ in the support is positive and the alphabet is finite. For rational masses this is the bounded-variable strong law of Chapter~4, because the symbols can be realised using finite digits. For arbitrary positive finite-alphabet masses, each symbol is obtained from a uniform by a finite partition of $[0,1)$; the bounded monotone step-function approximation proved after Proposition~\ref{prop:inverse} in Chapter~5 applies to its information content once the symbol probabilities are arranged in nondecreasing order. Thus the averages converge almost surely to $\mathbb E[\log_2(1/p_X)]=H(X)$ without silently assuming finite binary expansions for arbitrary probabilities.

\begin{theorem}[typical sequences]\label{thm:aep}
Let $\varepsilon>0$. The set of strings of length $n$ whose probability $p$ satisfies
\[
 2^{-n(H(X)+\varepsilon)}\le p\le 2^{-n(H(X)-\varepsilon)}
\]
has probability tending to $1$. Denote this set by $T_n(\varepsilon)$. For all sufficiently large $n$ its number of strings satisfies
\[
 \frac12\,2^{n(H(X)-\varepsilon)}<|T_n(\varepsilon)|\le2^{n(H(X)+\varepsilon)}.
\]
So nearly all the probability is carried by a set whose size grows exponentially at rate $H(X)$.
\end{theorem}
\begin{proof}
The information sum is an average of independent bounded variables. Its variance is $\operatorname{Var}(\log_2(1/p_X))/n$, so Chebyshev's inequality from Volume~1 proves directly that $P(T_n(\varepsilon))\to1$, without requiring a separate almost-sure-to-in-probability theorem. Each string in $T_n(\varepsilon)$ has mass at least $2^{-n(H+\varepsilon)}$, and the total mass is at most $1$, so $|T_n(\varepsilon)|\le2^{n(H+\varepsilon)}$. Conversely, each such string has mass at most $2^{-n(H-\varepsilon)}$, while their total mass eventually exceeds $1/2$. Therefore $|T_n(\varepsilon)|>\frac12\,2^{n(H-\varepsilon)}$.
\end{proof}

For a fair coin, every one of the $2^n$ binary strings has mass $2^{-n}$: the typical set can be the whole space. For a biased coin, most strings cluster near the expected frequency of heads, and the number of typical strings grows like $2^{nH}$ rather than $2^n$. The theorem supplies both an upper and a lower exponential bound, not the exact count of any one finite sample.

\section{Prefix codes}

A binary code assigns to each symbol $a$ a word $c(a)$ in the alphabet $\{0,1\}$. It is a \emph{prefix code} when no codeword is a prefix of another. A receiver can then stop at the end of each word without a separating mark.

\begin{theorem}[Kraft]\label{thm:kraft}
If $c$ is a prefix code with word lengths $\ell(a)$, then $\sum_a 2^{-\ell(a)}\le 1$. Conversely, any positive integers $\ell(a)$ with that sum at most $1$ are the lengths of some prefix code.
\end{theorem}
\begin{proof}
In the full binary tree, the codeword of length $\ell$ blocks $2^{-\ell}$ of the leaves at every deeper level. Prefix-freeness means the blocked sets are disjoint, so the sum is at most $1$. Conversely, sort the requested lengths as $\ell_1\le\cdots\le\ell_m$. Suppose the first $i-1$ codewords have been assigned, with none the prefix of another. Among the $2^{\ell_i}$ binary words of length $\ell_i$, the earlier words block exactly $\sum_{j<i}2^{\ell_i-\ell_j}$ descendants. Since
\[
\sum_{j\le i}2^{-\ell_j}\le\sum_{j\le m}2^{-\ell_j}\le1,
\]
there are at most $2^{\ell_i}-1$ blocked words, so an unblocked word exists. Choose the first one in lexicographic order. Because every earlier word has length at most $\ell_i$ and we avoided its descendants, the new word creates no prefix conflict. Continue through all $m$ symbols.
\end{proof}

\begin{theorem}[Shannon's code]\label{thm:shannon}
Let $\ell(a)=\lceil\log_2(1/p_a)\rceil$. Then $\sum 2^{-\ell(a)}\le 1$, so a prefix code of these lengths exists, and its expected length $L$ satisfies $H(X)\le L<H(X)+1$.
\end{theorem}
\begin{proof}
$\ell(a)\ge\log_2(1/p_a)$, so $2^{-\ell(a)}\le p_a$ and the sum is at most $1$. Also $\ell(a)<\log_2(1/p_a)+1$, so $L=\sum p_a\ell(a)<\sum p_a\log_2(1/p_a)+1=H(X)+1$. The lower bound here follows term by term from $\log_2(1/p_a)\le\ell(a)$. More generally, \emph{every} binary prefix code has $L\ge H(X)$, even if its lengths were not chosen by rounding. Indeed, let $Z=\sum_a2^{-\ell(a)}\le1$ by Kraft, and set $q_a=2^{-\ell(a)}/Z$. Gibbs' inequality yields
\[
 H(X)\le\sum_a p_a\log_2(1/q_a)
 =\sum_a p_a\ell(a)+\log_2 Z\le L.
\]
\end{proof}

The tree shows a prefix code that attains the entropy exactly.

\begin{figure}[ht]
\centering
\fig{ch07-prefixtree}
\caption{One complete binary prefix code. The four symbols have probabilities $1/2$, $1/4$, $1/8$ and $1/8$ and receive words $0$, $10$, $110$ and $111$. No endpoint of a word is above another endpoint. Kraft's sum is $1$, and the average codeword length equals the entropy $7/4$ bits.}
\label{fig:prefix}
\end{figure}

For example, a message consisting of $a$, $b$, $c$, $d$ is transmitted as $0\,10\,110\,111$. A decoder reads the tree from its root, emits a symbol at a leaf, and starts again. The expected length per symbol is
\[
L=\tfrac12\cdot1+\tfrac14\cdot2+\tfrac18\cdot3+\tfrac18\cdot3=\tfrac74.
\]
The entropy has exactly the same value. This is the special case where each mass is a negative power of $2$.

\section{Coding long blocks}

Coding one symbol at a time wastes up to one bit per symbol because lengths must be integers. Coding a block of $n$ independent symbols reduces that overhead to less than one bit per \emph{block}. A block is one symbol from the finite alphabet $A^n$. Its entropy is
\[
H(X_1,\ldots,X_n)=-\sum_{a_1,\ldots,a_n}\prod_{i=1}^np_{a_i}
  \log_2\Bigl(\prod_{i=1}^np_{a_i}\Bigr)
=nH(X),
\]
by expanding the logarithm and summing the other coordinates to $1$. Apply Shannon's code to the block law. There is a prefix code with average length $L_n$ satisfying
\[
nH(X)\le L_n<nH(X)+1,
\qquad H(X)\le\frac{L_n}{n}<H(X)+\frac1n.
\]
So the entropy is the best achievable average number of bits per original symbol, in the limit of coding long independent blocks with prefix codes. This statement is a source-coding theorem. It makes no claim about dependence or communication through a noisy channel.

The typical-set theorem tells us what the code is exploiting. Most of the probability is carried by roughly $2^{nH}$ likely blocks, rather than by all $|A|^n$ possible blocks. The integer-length argument shows how that saving becomes a real code.

Random walks, in the next chapter, ask a different question about repeated randomness: not how many typical strings there are, but whether a running sum will ever return.

\section{Exercises}

\begin{exercise}[$\star\star$]\label{ex:entropymax}
For a Bernoulli$(p)$ variable, $0<p<1$, show that $H(p)=-p\log_2 p-(1-p)\log_2(1-p)$ has a maximum at $p=1/2$, and that the maximum value is $1$.
\end{exercise}


\begin{exercise}[$\star$]\label{ex:prefixoptimal}
Let $P(a)=1/2$, $P(b)=1/4$ and $P(c)=P(d)=1/8$. Check that the words $0$, $10$, $110$, $111$ form a prefix code, calculate its expected length, and prove that no binary prefix code has smaller expected length.
\end{exercise}

\begin{exercise}[$\star\star$]\label{ex:entropychain}
For finite random variables $X$ and $Y$, define $H(Y\mid X)=\sum_xP(X=x)H(Y\mid X=x)$, omitting zero-probability terms. Starting from $P(X=x,Y=y)=P(X=x)P(Y=y\mid X=x)$, prove the chain rule $H(X,Y)=H(X)+H(Y\mid X)$. Calculate the two sides when $X$ is a fair bit and $Y=X$.
\end{exercise}

\begin{exercise}[$\star\star$]\label{ex:entropyproduct}
Show that if $X$ and $Y$ are independent finite-alphabet variables, then $H(X,Y)=H(X)+H(Y)$. Explain why the independence hypothesis cannot be dropped, using $Y=X$ with $X$ a fair bit.
\end{exercise}

\begin{exercise}[$\star$]\label{ex:ambiguouscode}
The three codewords $0$, $01$, $10$ are not prefix-free. Show that the bit string $010$ can be decoded as two different pairs of source symbols. Compare with Figure~\ref{fig:prefix}.
\end{exercise}

\begin{exercise}[$\star\star$]\label{ex:typicalcount}
Let $X$ be Bernoulli$(1/4)$ and let $T_n(\varepsilon)$ be the typical set from Theorem~\ref{thm:aep}. Verify that $P(T_n(\varepsilon))\to1$ by Chebyshev, identify $H(X)$ explicitly, and use the two bounds in the theorem to bound $\log_2|T_n(\varepsilon)|/n$ for large $n$. State what happens as $\varepsilon\downarrow0$.
\end{exercise}

\section{Hints}

\begin{hint}{ex:entropymax}
Differentiate. The second derivative of $-p\ln p-(1-p)\ln(1-p)$ is negative.
\end{hint}


\begin{hint}{ex:prefixoptimal}
No selected codeword is the beginning of another. Compare its average length with entropy, then use Kraft and Gibbs for any competitor.
\end{hint}
\begin{hint}{ex:entropychain}
Inside the double sum for $H(X,Y)$, write $-\log_2 p_{xy}=-\log_2 p_x-\log_2 p_{y\mid x}$.
\end{hint}
\begin{hint}{ex:entropyproduct}
For independent coordinates, $P(Y=y\mid X=x)=P(Y=y)$. The copying case has only two joint outcomes rather than four.
\end{hint}
\begin{hint}{ex:ambiguouscode}
Compare the split $0\,|\,10$ with another split using the codeword $01$.
\end{hint}
\begin{hint}{ex:typicalcount}
The information of one head is $\log_2 4$ and of one tail is $\log_2(4/3)$. The information sum has bounded variance.
\end{hint}

\section{Solutions}

\begin{solution}{ex:entropymax}
\idea The binary entropy is concave, and its derivative vanishes only at $1/2$.
\writeup Work in natural units and convert at the end. Let $h(p)=-p\ln p-(1-p)\ln(1-p)$. Then
\[
 h'(p)=-\ln p+\ln(1-p)=\ln\frac{1-p}{p},\qquad h''(p)=-\frac1p-\frac1{1-p}<0.
\]
So $h'$ vanishes only at $p=1/2$, and $h$ has a maximum there. Now $h(1/2)=\ln 2$, and dividing by $\ln 2$ converts to bits, giving $H(1/2)=1$. For every other $p\in(0,1)$, $H(p)<1$.
\end{solution}


\begin{solution}{ex:prefixoptimal}
\idea Draw the leaves in the binary tree, then compare length with the entropy lower bound.
\writeup No word in $\{0,10,110,111\}$ begins another word in that set, so decoding stops unambiguously at each leaf. The mean length is $L=(1/2)1+(1/4)2+(1/8)3+(1/8)3=7/4$. The entropy is $H=(1/2)\cdot1+(1/4)\cdot2+2(1/8)\cdot3=7/4$. For any other prefix code with lengths $\ell_i$, Kraft gives $Z=\sum_i2^{-\ell_i}\le1$. Gibbs applied to $q_i=2^{-\ell_i}/Z$ yields $H\le L+\log_2 Z\le L$. Thus no other prefix code has smaller mean.
\end{solution}

\begin{solution}{ex:entropychain}
\idea Split the information of the pair into the information needed to identify $X$ and the information still needed for $Y$ after that.
\writeup Write $p_{xy}=p_xp_{y\mid x}$ whenever $p_{xy}>0$. Then, with zero-mass terms omitted,
\[
\begin{aligned}
H(X,Y)&=-\sum_{x,y}p_{xy}\log_2(p_xp_{y\mid x})\\
&=-\sum_xp_x\log_2p_x\sum_y p_{y\mid x}
  -\sum_xp_x\sum_y p_{y\mid x}\log_2p_{y\mid x}\\
&=H(X)+H(Y\mid X).
\end{aligned}
\]
When $X$ is a fair bit and $Y=X$, the pair has only $(0,0)$ and $(1,1)$, each of probability $1/2$, so $H(X,Y)=1$. Knowing $X$ determines $Y$, so $H(Y\mid X)=0$ and $H(X)=1$.
\end{solution}

\begin{solution}{ex:entropyproduct}
\idea Independence says that knowing $X$ does not change the law of $Y$.
\writeup If $X$ and $Y$ are independent, then for every $x$ of positive probability the conditional masses $p_{y\mid x}=p_y$, hence $H(Y\mid X)=H(Y)$. The chain rule of Exercise~\ref{ex:entropychain} gives $H(X,Y)=H(X)+H(Y)$. If $Y=X$ and $X$ is fair, both marginal entropies are $1$ bit but the joint entropy is only $1$ bit. The extra bit of marginal entropy is not new information.
\end{solution}

\begin{solution}{ex:ambiguouscode}
\idea A prefix conflict can create a second parsing of the same transmitted bits.
\writeup The bit string $010$ is either $0\,|\,10$ (the first and third symbols) or $01\,|\,0$ (the second and first symbols). Both concatenations produce the same bits. In contrast, no codeword in Figure~\ref{fig:prefix} is a prefix of another, and the decoder knows when each word ends.
\end{solution}

\begin{solution}{ex:typicalcount}
\idea Information per symbol is a bounded random variable; apply Chebyshev and then count the total mass of the typical set.
\writeup Put $I_i=\log_2(1/p_{X_i})$, so $I_i$ equals $2$ with probability $1/4$ and $\log_2(4/3)$ with probability $3/4$. Its mean is
\[
H(X)=\frac14\cdot2+\frac34\log_2(4/3)
= -\frac14\log_2(1/4)-\frac34\log_2(3/4).
\]
Its variance is finite. Since the $I_i$ are independent, Chebyshev gives
$P(|n^{-1}\sum_{i=1}^n I_i-H(X)|>\varepsilon)\le
\operatorname{Var}(I_1)/(n\varepsilon^2)\to0$.
Therefore $P(T_n(\varepsilon))\to1$. For all sufficiently large $n$ the typical-set count gives
\[
H(X)-\varepsilon-\frac1n
<\frac{\log_2|T_n(\varepsilon)|}{n}
\le H(X)+\varepsilon,
\]
where $\frac12=2^{-1}$ supplies the $1/n$. First let $n\to\infty$ and then $\varepsilon\downarrow0$: the exponential growth rate is $H(X)$.
\end{solution}

\section*{Sources and further reading}

\begin{enumerate}
\item C.~E.~Shannon, A mathematical theory of communication, \emph{Bell System Tech.\ J.} \textbf{27} (1948), 379--423 and 623--656. The entropy, the source-coding bound, and typical sequences for a finite alphabet.
\end{enumerate}
